%---
% slug: terel
% title: Temporal Regularized Learning
% category: research
% eyebrow: LOCAL LEARNING
% summary: TeReL turns a soft-SFA target residual into a neuron state whose endpoint-local factors update feedforward and lateral synapses without errors crossing layers or time.
% draft: false
% date: 2026-07-15
% rank: 2
%---
\documentclass[11pt]{article}
\usepackage{publication-web}
\begin{document}

\publicationtitle{Temporal Regularized Learning: Deep Slow-Feature Learning Local in Space and Time}
\publicationauthors{Davide R. Wiest}
\publicationaffiliations{Axym Labs; Darmstadt, Germany}

\begin{publicationhead}
\makepublicationtitle
\end{publicationhead}

\publicationfigure[\linewidth]
  {figures/terel-lead.png}{/work/terel-lead.png}
  {Three regular stages show a soft-SFA target becoming a neuron state and then two local synaptic updates.}
  {TeReL converts soft-SFA forces into neuron-local synaptic signals. A detached regularized target defines a preactivation neuron state; one inhibitory pass modifies that state, which then supplies both the feedforward and anti-Hebbian lateral changes.}

\publicationabstract{Slow Feature Analysis learns persistent factors by making adjacent representations similar while preserving variance and decorrelating coordinates. Temporal Regularized Learning (TeReL) derives a detached target from these three forces and maps its residual to the preactivation of each neuron. The outer product of this neuron state and presynaptic activity is the exact feedforward gradient without lateral coupling. One learned inhibitory pass modifies the state, while pairs of the same states drive an anti-Hebbian lateral rule. TeReL processes one observation at a time, retains fixed detached state, uses plain stochastic gradient descent without optimizer moments, and sends no error between learned layers.}

\publicationactions{%
  \publicationaction{/files/terel.pdf}{Read paper}
  \publicationaction{https://github.com/Axym-Labs/TeReL}{View code}
  \publicationaction{https://doi.org/10.5281/zenodo.18673107}{Open DOI}%
}

\section{Method}

For one observation, temporal similarity, variance preservation, and
decorrelation define a regularized activation target. The residual to this
target becomes a preactivation neuron state. Multiplying that postsynaptic
state by presynaptic activity gives the feedforward update, while products of
two neuron states train a same-layer inhibitory operator. Spatial locality
means that each synapse reads quantities available at its endpoints, possibly
after same-layer dynamics; temporal locality means that the update reads fixed
detached state rather than an activation graph extending into the past.

\section{Main results}

On a label-ordered MNIST stream, samplewise TeReL reaches
$95.84\pm0.07\%$ held-out accuracy after two presentations. On a train-derived
validation split, one residual-state inhibitory pass improves an otherwise
matched samplewise rule by 1.42 percentage points without sacrificing
effective rank. The less constrained minibatch reference, TeReL-Offline,
reaches $97.30\pm0.07\%$, compared with $98.34\pm0.08\%$ for matched
backpropagation and $96.98\pm0.10\%$ for Local SupCon. Objective ablations and
a direct-covariance control identify the roles of the three soft-SFA forces and
the lagged decorrelation signal.

\section{Citation}

\begin{verbatim}
@misc{wiest2026terel,
  author = {Wiest, Davide R.},
  title  = {Temporal Regularized Learning: Deep Slow-Feature
            Learning Local in Space and Time},
  year   = {2026},
  doi    = {10.5281/zenodo.18673107},
  url    = {https://doi.org/10.5281/zenodo.18673107}
}
\end{verbatim}

\end{document}
